% 第 2 章 随机化实验
\providecommand{\ciLeafA}[6]{%
  \begin{scope}[shift={(#1,#2)}]
    \draw[thin] (0,0) ellipse (1.30 and 0.62);
    \draw[thin] (-1.26,-0.15) -- (1.26,0.15);
    \node[font=\footnotesize,inner sep=0pt] at (0,0.31) {#3};
    \node[font=\footnotesize,inner sep=0pt] at (0,-0.31) {#4};
    \node[font=\footnotesize,anchor=west,inner sep=0pt] at (1.40,0.31) {#5};
    \node[font=\footnotesize,anchor=west,inner sep=0pt] at (1.40,-0.31) {#6};
  \end{scope}%
}
\providecommand{\ciLeafB}[6]{%
  \begin{scope}[shift={(#1,#2)}]
    \draw[thin] (0,0) ellipse (1.00 and 0.44);
    \draw[thin] (-0.97,-0.11) -- (0.97,0.11);
    \node[font=\scriptsize,inner sep=0pt] at (0,0.21) {#3};
    \node[font=\scriptsize,inner sep=0pt] at (0,-0.21) {#4};
    \node[font=\scriptsize,anchor=west,inner sep=0pt] at (1.08,0.21) {#5};
    \node[font=\scriptsize,anchor=west,inner sep=0pt] at (1.08,-0.21) {#6};
  \end{scope}%
}

\chapter{随机化实验}

你抬头看天，身边的行人会不会也跟着抬头？这个问题包含了任何因果问题的主要构件：我们想知道，某个行动（你抬头）是否影响某个结局（别人抬头），并且是在一个特定的总体（比如 2019 年马德里的居民）之中。假设我们请你设计一项科学研究来回答这个问题。你想了想，说：“这不算什么挑战。我可以站在人行道上，每逢有人走近就掷一次硬币。正面，我就抬头；反面，我就平视前方。这样重复几千次。如果我抬头之后 10 秒内抬头的行人比例，高于我没有抬头时抬头的行人比例，我就可以断定：我的抬头对别人抬头有因果效应。对了，我大概还得雇个助手，在我抬头的时候记录路人的反应。”研究做完，你发现：你抬头时，55\% 的行人跟着抬头；你平视前方时，只有 1\% 的行人抬头。

你给出的方案，是做一项\emph{随机化实验}。说它是实验，是因为所关心的行动（抬头）由研究者（你）亲自实施；说它是随机化的，是因为对每个研究对象（行人）是否采取行动，由一个随机装置（掷硬币）决定。并非所有实验都是随机化的。比如，你可以在男人走近时抬头、女人走近时平视前方。这样，行动的分配遵循的是一条确定性规则（男人抬头、女人平视），而不是随机机制。但如果你做的是非随机化实验，你的发现就远不会这么有说服力。一旦行动由行人的性别决定，批评者就可以争辩说，男人和女人的“抬头”行为本就不同（你抬头之后，女人抬头的频率也许低于男人），于是你的研究实质上比较了两个“不可比”的人群。本章要说明的，正是随机化为何能带来令人信服的因果推断。

\section{随机化}

在真实世界的研究中，我们不会同时知道宙斯的两个潜在结局——治疗之下的 $Y^{a=1}$ 与未治疗之下的 $Y^{a=0}$。我们能知道的，只是他在实际接受的治疗取值 $A$ 之下的观测结局 $Y$。表~2.1 汇总了我们这个 20 人总体所能获得的信息。每个人两个反事实结局中只有一个已知：与他实际接受的治疗水平相对应的那一个；另一个反事实结局的数据缺失。如上一章所述，这种数据缺失带来了麻烦，因为计算效应度量似乎需要两个反事实结局的取值。表~2.1 的数据只够用来计算关联度量。

\mnote{Neyman (1923) 把反事实理论用于借助随机化实验估计因果效应。}

\begin{table}[t]
\centering
\caption{心脏移植研究的观测数据（问号表示该反事实结局未被观测）}
\label{tab:2.1}
\small
\begin{tabular}{lcccc}
\toprule
& $A$ & $Y$ & $Y^{a=0}$ & $Y^{a=1}$\\
\midrule
Rheia      & 0 & 0 & 0 & ?\\
Kronos     & 0 & 1 & 1 & ?\\
Demeter    & 0 & 0 & 0 & ?\\
Hades      & 0 & 0 & 0 & ?\\
Hestia     & 1 & 0 & ? & 0\\
Poseidon   & 1 & 0 & ? & 0\\
Hera       & 1 & 0 & ? & 0\\
Zeus       & 1 & 1 & ? & 1\\
Artemis    & 0 & 1 & 1 & ?\\
Apollo     & 0 & 1 & 1 & ?\\
Leto       & 0 & 0 & 0 & ?\\
Ares       & 1 & 1 & ? & 1\\
Athena     & 1 & 1 & ? & 1\\
Hephaestus & 1 & 1 & ? & 1\\
Aphrodite  & 1 & 1 & ? & 1\\
Polyphemus & 1 & 1 & ? & 1\\
Persephone & 1 & 1 & ? & 1\\
Hermes     & 1 & 0 & ? & 0\\
Hebe       & 1 & 0 & ? & 0\\
Dionysus   & 1 & 0 & ? & 0\\
\bottomrule
\end{tabular}
\end{table}

随机化实验与任何其他真实世界研究一样，产生的数据中反事实结局是缺失的，如表~2.1 所示。然而，随机化保证这些缺失是偶然发生的。因此，尽管数据缺失，随机化实验中仍可计算效应度量——更严格地说，仍可一致地估计效应度量。下面我们把这一点讲得更精确些。

设想图~1.1 中菱形所代表的总体近乎无限大，我们为其中每个人掷一次硬币：反面归入白色组，正面归入灰色组。注意这不是一枚公平的硬币，正面的概率不到 50\%，所以落到灰色组的人比白色组少。接着，我们让研究助手给白色组的人施用所关心的治疗（$A=1$），给灰色组的人施用安慰剂（$A=0$）。五天后研究结束时，我们计算各组的死亡风险：$\Pr[Y=1\given A=1]=0.3$，$\Pr[Y=1\given A=0]=0.6$。关联风险比为 $0.3/0.6=0.5$，关联风险差为 $0.3-0.6=-0.3$。我们假定这是一项在其他各方面都理想的随机化实验：无失访，研究期间完全依从所分配的治疗，治疗只有一种版本，且分配是双盲的（见第 9 章）。理想的随机化实验并不现实，但用来引入因果推断的若干关键概念很有用。本书后面再讨论更接近现实的随机化实验。

现在设想，假如研究助手误解了我们的指示，给灰色组而不是白色组施用了治疗，会怎样？假定研究结束后我们才发现这一误会。治疗状态的这种颠倒，会影响我们的结论吗？丝毫不会。我们仍会发现，受治疗者（此时是灰色组）的风险 $\Pr[Y=1\given A=1]$ 为 0.3，未受治疗者（此时是白色组）的风险 $\Pr[Y=1\given A=0]$ 为 0.6，关联度量不变。

由于个体是被随机分配到白色组与灰色组的，无论接受治疗的是白色组、接受安慰剂的是灰色组，还是反过来，暴露者中的死亡比例 $\Pr[Y=1\given A=1]$ 都应当相同。当组别归属是随机化的，究竟哪一组接受了治疗，对 $\Pr[Y=1\given A=1]$ 的取值无关紧要。同样的道理当然也适用于 $\Pr[Y=1\given A=0]$。正式地说，我们称这两组是\emph{可交换的}。可交换性意味着：假如白色组的人接受了分配给灰色组的那种治疗，白色组的死亡风险本会与灰色组的死亡风险相同。也就是说，对 $a=0$ 与 $a=1$ 都有：受治疗者在潜在治疗取值 $a$ 之下的风险 $\Pr[Y^{a}=1\given A=1]$，等于未受治疗者在潜在治疗取值 $a$ 之下的风险 $\Pr[Y^{a}=1\given A=0]$。这些（条件）风险在总体中由治疗状态划分的各个子集里都相等，一个显然的推论是：它们必等于整个总体在治疗取值 $a$ 之下的（边际）风险，即 $\Pr[Y^{a}=1\given A=1]=\Pr[Y^{a}=1\given A=0]=\Pr[Y^{a}=1]$。既然治疗取值 $a$ 之下的反事实风险在 $A=1$ 与 $A=0$ 两组相同，我们就说，实际治疗 $A$ 不能预测反事实结局 $Y^{a}$。等价地，可交换性意味着反事实结局与实际治疗独立，即对所有 $a$ 都有 $Y^{a}\indep A$。随机化之所以备受重视，正是因为它有望产生可交换性。当受治疗者与未受治疗者可交换时，我们有时说治疗是\emph{外生的}（exogenous），于是外生性常被用作可交换性的同义词。

\mnote{可交换性：对所有 $a$ 都有 $Y^{a}\indep A$。其他版本见技术注记 2.1。}

\renewcommand{\TPname}{技术注记（Technical Point）}
\begin{techpoint}[完整可交换性与均值可交换性]
随机化使 $Y^{a}$ 联合独立于 $A$，这蕴含（但不被蕴含于）对每个 $a$ 的可交换性 $Y^{a}\indep A$。正式地，令 $A=\{a,a',a'',\dots\}$ 表示总体中出现的全部治疗取值之集，$Y^{A}=\{Y^{a},Y^{a'},Y^{a''},\dots\}$ 表示全部反事实结局之集。随机化使 $Y^{A}\indep A$。我们把这种联合独立性称为\emph{完整可交换性}。对二值治疗，$A=\{0,1\}$，完整可交换性即 $(Y^{a=1},Y^{a=0})\indep A$。

对二值结局与二值治疗，可交换性 $Y^{a}\indep A$ 也可写作 $\Pr[Y^{a}=1\given A=1]=\Pr[Y^{a}=1\given A=0]$，或等价地写作 $\E[Y^{a}\given A=1]=\E[Y^{a}\given A=0]$（对所有 $a$）。我们把最后一个等式称为\emph{均值可交换性}。对连续结局，可交换性 $Y^{a}\indep A$ 蕴含均值可交换性 $\E[Y^{a}\given A=a']=\E[Y^{a}]$，但均值可交换性不蕴含可交换性，因为均值以外的分布参数（如方差）未必独立于治疗。

要证明 $\E[Y^{a}]=\E[Y\given A=a]$，既不需要完整可交换性 $Y^{A}\indep A$，也不需要可交换性 $Y^{a}\indep A$，均值可交换性已经足够。如正文所述，证明分两步。第一步，由一致性有 $\E[Y\given A=a]=\E[Y^{a}\given A=a]$。第二步，由均值可交换性有 $\E[Y^{a}\given A=a]=\E[Y^{a}]$。由于本章所用的二值结局下，可交换性与均值可交换性是同一个概念，我们通篇采用较短的名称“可交换性”。
\end{techpoint}
\renewcommand{\TPname}{技术注记}

上一段论证的是：在可交换性之下，总体中白色那部分的治疗之下的反事实风险，等于整个总体的治疗之下的反事实风险。可是白色组在治疗之下的风险根本算不上反事实——白色组确实接受了治疗！因此，我们的理想随机化实验使我们能够算出总体在治疗之下的反事实风险 $\Pr[Y^{a=1}=1]$，因为它等于受治疗者的风险 $\Pr[Y=1\given A=1]=0.3$。也就是说，受治疗者（菱形中白色的部分）的风险，与所有人都受治疗（从而整个菱形都变白）时的风险相同。同样的道理当然适用于未受治疗者：总体在未治疗之下的反事实风险 $\Pr[Y^{a=0}=1]$ 等于未受治疗者的风险 $\Pr[Y=1\given A=0]=0.6$。因果风险比为 0.5，因果风险差为 $-0.3$。在理想的随机化实验中，关联即因果。

随机化实验中 $Y^{a}\indep A$ 还有另一种解释。反事实结局 $Y^{a}$ 就像一个人的基因组成，可以看作此人在治疗 $A$ 被随机分配之前就已具备的固定特征。这是因为 $Y^{a}$ 编码的是“若被分配到治疗 $a$，此人的结局本会是什么”，因而并不依赖于你后来实际接受的治疗。由于治疗 $A$ 是随机化的，它既独立于你的基因，也独立于 $Y^{a}$。$Y^{a}$ 与基因组成的不同之处在于：即便只在概念上，你也只能在治疗给出之后、并且只有当此人的治疗 $A$ 等于 $a$ 时，才能得知 $Y^{a}$ 的取值。

往下读之前，请务必弄清 $Y^{a}\indep A$ 与 $Y\indep A$ 的区别。\mnote{注意：$Y^{a}\indep A$ 与 $Y\indep A$ 不同。}可交换性 $Y^{a}\indep A$ 的定义是反事实结局与观测治疗之间的独立性。再一次说明：这意味着，如果受治疗者与未受治疗者接受相同的治疗水平（$a=0$ 或 $a=1$），他们本会经历相同的死亡风险。但是，反事实结局与观测治疗之间的独立性 $Y^{a}\indep A$，并不意味着观测结局与观测治疗之间的独立性 $Y\indep A$。例如，在一项可交换性 $Y^{a}\indep A$ 成立、且治疗对结局有因果效应的随机化实验中，$Y\indep A$ 并不成立，因为治疗与观测结局有关联。

\mnote{设某些个体的处理有因果效应，即 $Y^{a=1}\neq Y^{a=0}$。由 $Y=Y^{A}$，把 $Y^{a}$ 中的 $a$ 取为观测到的治疗 $A$，即得观测到的 $Y^{A}$；它依赖于 $A$，因而不会与 $A$ 独立。}

\renewcommand{\FPname}{精要（Fine Point）}
\begin{finepoint}[交叉实验]
设想我们要估计宙斯使用闪电 $A$ 对其血压 $Y$ 的个体因果效应。定义反事实结局 $Y^{a=1}$ 与 $Y^{a=0}$：若宙斯召唤（或不召唤）闪电之后血压暂时升高，则分别取 1。假定我们说服宙斯，只在我们的建议下使用闪电。昨天早上，我们让宙斯召唤闪电（$a=1$），之后他的血压升高了。今天早上，我们让宙斯不要使用闪电（$a=0$），他的血压没有升高。我们做的是一项\emph{交叉实验}：同一个体的结局在两个治疗取值下被先后观测。有人可能争辩说，既然我们观测到了宙斯的两个反事实结局 $Y^{a=1}=1$ 与 $Y^{a=0}=0$，使用闪电对宙斯的血压就有因果效应。然而，除非下一段给出的三条很强的假定 i)–iii) 成立，这一论证一般是错误的。

在交叉实验中，个体在两个或多个时段被观测，比如 $t=0$ 与 $t=1$。个体 $i$ 在每个时段 $t$ 接受不同的治疗取值 $A_{it}$。记 $Y_{i1}^{a_0a_1}$ 为个体 $i$ 在 $t=1$ 时的（确定性）反事实结局：若在 $t=1$ 接受治疗 $a_1$、在 $t=0$ 接受治疗 $a_0$。$t=0$ 时的 $Y_{i0}^{a_0}$ 类似定义。若下列三个条件成立，个体因果效应 $Y_{it}^{a_t=1}-Y_{it}^{a_t=0}$ 可识别：i) 无治疗的延续效应（carryover effect）：$Y_{i1}^{a_0a_1}=Y_{i1}^{a_1}$；ii) 个体因果效应不依赖于时间：$Y_{it}^{a_t=1}-Y_{it}^{a_t=0}=\alpha_i$，$t=0,1$；iii) 未治疗之下的反事实结局不依赖于时间：$Y_{it}^{a_t=0}=\beta_i$，$t=0,1$。在这些条件之下，若个体在 $t=1$ 受治疗（$A_{i1}=1$）而在 $t=0$ 未受治疗（$A_{i0}=0$），则由一致性，$Y_{i1}-Y_{i0}$ 就是个体因果效应，因为
\[
Y_{i1}-Y_{i0}=Y_{i1}^{a_1=1}-Y_{i0}^{a_0=0}
=Y_{i1}^{a_1=1}-Y_{i1}^{a_1=0}+Y_{i1}^{a_1=0}-Y_{i0}^{a_0=0}
=\alpha_i+\beta_i-\beta_i=\alpha_i.
\]
同理，若 $A_{i1}=0$、$A_{i0}=1$，则 $Y_{i0}-Y_{i1}=\alpha_i$ 就是个体层面的因果效应。

条件 (i) 意味着结局 $Y_{it}^{a_t}$ 骤然出现，并在下一个时段之前完全消退。因此，交叉实验不能用来研究心脏移植这一不可逆行动对死亡这一不可逆结局的效应。另见精要 3.2。
\end{finepoint}
\renewcommand{\FPname}{精要}

表~2.1 的心脏移植研究中，可交换性成立吗？要回答这个问题，需要检查 $a=0$ 与 $a=1$ 时 $Y^{a}\indep A$ 是否成立。先看 $a=0$。假定表~1.1 的反事实数据可以拿到，我们就能算出：13 名受治疗者在未治疗之下的死亡风险为 $\Pr[Y^{a=0}=1\given A=1]=7/13$，7 名未受治疗者在未治疗之下的死亡风险为 $\Pr[Y^{a=0}=1\given A=0]=3/7$。未治疗之下的死亡风险在受治疗者中高于未受治疗者，即 $7/13>3/7$，可见受治疗者的预后比未受治疗者更差，也就是说，受治疗者与未受治疗者不可交换。从数学上讲，我们已经证明可交换性 $Y^{a}\indep A$ 在 $a=0$ 时不成立。（你可以自行验证 $a=1$ 时它也不成立。）因此，本段开头那个问题的答案是“不成立”。

可是在真实世界里，能拿到的只有表~2.1 的观测数据，而不是表~1.1 的反事实数据。既然表~2.1 不足以计算诸如受治疗者在未治疗之下的风险 $\Pr[Y^{a=0}=1\given A=1]$ 这类反事实风险，我们一般无法判断自己的研究中可交换性是否成立。不过，暂且假定我们真的能拿到表~1.1，并判定这项心脏移植研究中可交换性不成立。那么，能否由此断定我们的研究不是随机化实验？不能，理由有二。第一，你大概已经想到了：一项 20 人的研究太小，得不出确定的结论。抽样变异引起的随机波动几乎可以解释任何现象。随机变异留到第 10 章讨论；在此之前，我们假定总体中每个人都代表 10 亿个与他完全相同的人。第二，即便在无限样本中可交换性不成立，一项研究仍可能是随机化实验。只是与本节描述的那类随机化实验不同，它必须是研究者用不止一枚硬币来随机分配治疗的随机化实验。下一节就讨论用不止一枚硬币的随机化实验。

\mnote{提请注意：我们对随机化实验的讨论针对的是总体或平均因果效应，因为个体因果效应一般不可识别。见精要 2.1。}

\section{条件随机化}

表~2.2 给出我们这项心脏移植随机化研究的数据。除治疗 $A$（接受移植取 1，否则取 0）与结局 $Y$（死亡取 1，否则取 0）之外，表~2.2 还包含预后因子 $L$（病情危重取 1，否则取 0）的数据，$L$ 是在分配治疗之前测得的。下面考虑两种互斥的研究设计，并讨论表~2.2 的数据是否可能来自其中任何一种。

\begin{table}[t]
\centering
\caption{心脏移植随机化研究的数据（含治疗前的预后因子 $L$）}
\label{tab:2.2}
\small
\begin{tabular}{lccc}
\toprule
& $L$ & $A$ & $Y$\\
\midrule
Rheia      & 0 & 0 & 0\\
Kronos     & 0 & 0 & 1\\
Demeter    & 0 & 0 & 0\\
Hades      & 0 & 0 & 0\\
Hestia     & 0 & 1 & 0\\
Poseidon   & 0 & 1 & 0\\
Hera       & 0 & 1 & 0\\
Zeus       & 0 & 1 & 1\\
Artemis    & 1 & 1 & 0\\
Apollo     & 1 & 0 & 1\\
Leto       & 1 & 1 & 0\\
Ares       & 1 & 1 & 1\\
Athena     & 1 & 1 & 1\\
Hephaestus & 1 & 1 & 1\\
Aphrodite  & 1 & 1 & 1\\
Polyphemus & 1 & 1 & 1\\
Persephone & 1 & 1 & 1\\
Hermes     & 1 & 0 & 1\\
Hebe       & 1 & 0 & 0\\
Dionysus   & 1 & 1 & 0\\
\bottomrule
\end{tabular}
\end{table}

在设计 1 中，我们从总体中随机抽取 65\% 的个体，给每个被抽中的人移植一颗新心脏。这可以解释为什么 20 人中有 13 人受治疗。在设计 2 中，我们先把所有人分为病情危重（$L=1$）与病情不危重（$L=0$）两类，然后从危重者中随机抽取 75\%、从不危重者中随机抽取 50\%，给每个被抽中的人移植一颗新心脏。这可以解释为什么危重的 12 人中有 9 人、不危重的 8 人中有 4 人受治疗。

两种设计都是随机化实验。设计 1 正是 2.1 节描述的那类随机化实验。在这种设计下，我们用一枚硬币为所有个体分配治疗（例如反面受治疗、正面不受治疗）：一枚偏斜的硬币，反面的概率为 0.65，于是 65\% 的个体受治疗。在设计 2 下，我们不为所有人用同一枚硬币，而是对危重者用一枚反面概率 0.75 的硬币，对不危重者用另一枚反面概率 0.50 的硬币。我们把设计 2 这类实验称为\emph{条件随机化实验}，因为我们所用的若干个随机化概率依赖于（以之为条件）变量 $L$ 的取值。我们把设计 1 这类实验称为\emph{边际随机化实验}，因为我们只用单一的无条件（边际）随机化概率，它对所有个体都相同。

如上一节所述，边际随机化实验有望带来受治疗者与未受治疗者的可交换性：
\[
\Pr[Y^{a}=1\given A=1]=\Pr[Y^{a}=1\given A=0]\quad\text{或}\quad Y^{a}\indep A\ \text{（对所有 }a\text{）}.
\]
相反，条件随机化实验一般不会带来受治疗者与未受治疗者的可交换性，因为按设计，两组中预后不良者的比例可能不同。表~2.2 的数据不可能来自边际随机化实验，因为受治疗者中有 69\%、未受治疗者中有 43\% 病情危重。这种不平衡表明，受治疗者若一直未受治疗，其死亡风险本会高于未受治疗者的死亡风险。也就是说，治疗 $A$ 能预测未治疗之下的反事实死亡风险，可交换性 $Y^{a}\indep A$ 不成立。既然我们的研究是一项随机化实验，你就可以有把握地断定：它是一项以 $L$ 为条件进行随机化的随机化实验。

我们的条件随机化实验，不过是两项彼此独立的边际随机化实验的组合：一项在病情危重（$L=1$）的个体子集中进行，另一项在病情不危重（$L=0$）的个体子集中进行。先看危重者子集中进行的随机化实验。在这个子集里，受治疗者与未受治疗者可交换。正式地说，给定他们在分配治疗时都病情危重，每个治疗取值 $a$ 之下的反事实死亡风险在受治疗者与未受治疗者中相同，即
\begin{gather*}
\Pr[Y^{a}=1\given A=1,L=1]=\Pr[Y^{a}=1\given A=0,L=1]\\
\text{或}\quad Y^{a}\indep A\given L=1\quad\text{（对所有 }a\text{）},
\end{gather*}
其中 $Y^{a}\indep A\given L=1$ 表示给定 $L=1$ 时 $Y^{a}$ 与 $A$ 独立。同样，随机化也保证受治疗者与未受治疗者在病情不危重的个体子集中可交换，即 $Y^{a}\indep A\given L=0$。当 $Y^{a}\indep A\given L=l$ 对所有取值 $l$ 都成立时，我们简写为 $Y^{a}\indep A\given L$。\mnote{条件可交换性：对所有 $a$ 都有 $Y^{a}\indep A\given L$。}可见，条件随机化虽不保证无条件（或边际）可交换性 $Y^{a}\indep A$，却保证了在变量 $L$ 各水平之内的\emph{条件可交换性} $Y^{a}\indep A\given L$。总之，边际随机化（设计 1）同时产生边际可交换性与条件可交换性，而条件随机化（设计 2）只产生条件可交换性。

\mnote{若 $A=1$，则 $Y^{a=0}$ 缺失；若 $A=0$，则 $Y^{a=1}$ 缺失。若 $\Pr[A=a\given L,Y^{a=1},Y^{a=0}]=\Pr[A=a]$，称数据\emph{完全随机缺失}（MCAR），这在边际随机化实验中成立。若给定完整数据 $(L,Y^{a=1},Y^{a=0})$ 时 $A=a$ 的概率只依赖于 $A=a$ 时本可观测到的数据 $(L,Y^{a})$，称数据\emph{随机缺失}（MAR）。事实上，MAR 蕴含 $\Pr[A=a\given L,Y^{a=1},Y^{a=0}]=\Pr[A=a\given L]$，这在条件随机化实验中成立：由 MAR，$\Pr[A=1\given L,Y^{a=1},Y^{a=0}]$ 不可能依赖于 $Y^{a=0}$，而 $1-\Pr[A=1\given L,Y^{a=1},Y^{a=0}]=\Pr[A=0\given L,Y^{a=1},Y^{a=0}]$ 不可能依赖于 $Y^{a=1}$。MCAR、MAR 与 MNAR（非随机缺失）这三个术语由 Rubin (1976) 以及 Marini、Olsen 与 Rubin (1980) 引入。}

边际可交换性之下如何计算效应度量，我们已经知道：在边际随机化实验中，因果风险比 $\Pr[Y^{a=1}=1]/\Pr[Y^{a=0}=1]$ 等于关联风险比 $\Pr[Y=1\given A=1]/\Pr[Y=1\given A=0]$，因为可交换性保证治疗水平 $a$ 之下的反事实风险 $\Pr[Y^{a}=1]$ 等于接受治疗水平 $a$ 者的观测风险 $\Pr[Y=1\given A=a]$。因此，假如表~2.2 的数据是在边际随机化实验中收集的，因果风险比便可由 $A$ 与 $Y$ 的数据直接算出：$(7/13)/(3/7)=1.26$。问题在于，条件随机化实验中如何计算因果风险比。请记住，条件随机化实验不过是在总体不同子集（$L=1$ 与 $L=0$）中分别进行的两项（或更多）边际随机化实验的组合。因此我们有两种选择。

第一种，计算总体各个子集（层）中的平均因果效应。既然在每个子集内关联即因果，危重者的层特异因果风险比 $\Pr[Y^{a=1}=1\given L=1]/\Pr[Y^{a=0}=1\given L=1]$ 就等于危重者的层特异关联风险比 $\Pr[Y=1\given L=1,A=1]/\Pr[Y=1\given L=1,A=0]$；$L=0$ 时同理。我们把这种计算层特异因果效应的方法称为\emph{分层}。注意，$L=1$ 子集的层特异因果风险比可能与 $L=0$ 的因果风险比不同。此时我们说，治疗的效应被 $L$ 修饰，或称存在 $L$ 的\emph{效应修饰}，或称治疗效应在 $L$ 各水平间存在\emph{异质性}。分层与效应修饰在第 4 章详述。

第二种，计算整个总体中的平均因果效应 $\Pr[Y^{a=1}=1]/\Pr[Y^{a=0}=1]$，这也是我们迄今一直在做的。我们主要关心层特异平均因果效应，还是整个总体的平均因果效应，取决于第 4 章与第三部分详加讨论的实际与理论考量。举例来说，如果你预计无法获得未来个体的 $L$ 信息（例如 $L$ 测量起来很贵），从而你的治疗决策不可能依赖于 $L$ 的取值，那么你关心的可能是整个总体的平均因果效应，而不是层特异平均因果效应。在第 4 章之前，我们只关注整个总体的平均因果效应。接下来两节说明如何用条件随机化实验的数据计算整个总体的平均因果效应。关于风险期的讨论另见精要 2.2。

\begin{finepoint}[风险期]
我们把风险定义为特定时期内发生所关心结局的个体所占的比例。例如，受治疗者的 5 日死亡风险 $\Pr[Y=1\given A=1]$，就是受治疗者中在随访头五天内死亡者所占的比例。本书通篇常常在风险首次定义时指明时期（如 5 日），此后为简洁起见略去。也就是说，我们可能只说“死亡风险”，而不说“5 日死亡风险”。

下面这个例子说明指明风险期的重要性。假设为量化抗生素治疗对感染鼠疫杆菌的老年人死亡率的因果效应，做了一项随机化实验。一位研究者分析数据后得出结论：因果风险比为 0.05，即平均而言抗生素使死亡率降低 95\%。另一位研究者也分析了同一批数据，却得出结论：因果风险比为 1，即抗生素对死亡率没有平均因果效应。两位研究者都对。第一位算的是 1 年风险之比，第二位算的是 100 年风险之比。无论是否接受治疗，100 年风险当然都是 1。当我们说某种治疗对死亡率有因果效应时，我们的意思是治疗推迟了死亡，而不是阻止了死亡。
\end{finepoint}

\section{标准化}

我们的心脏移植研究是一项条件随机化实验：研究者用随机程序分配心脏（$A=1$），对病情不危重（$L=0$）的 8 人分配概率为 50\%，对病情危重（$L=1$）的 12 人分配概率为 75\%。先看病情不危重的那 8 个人——请记住，他们实际上是 80 亿人的平均代表。这一组中，受治疗者的死亡风险为 $\Pr[Y=1\given L=0,A=1]=1/4$，未受治疗者的死亡风险为 $\Pr[Y=1\given L=0,A=0]=1/4$。由于 $L=0$ 组的人是随机分配治疗的，即 $Y^{a}\indep A\given L=0$，观测风险等于反事实风险。也就是说，在 $L=0$ 组中，受治疗者的风险等于所有人都受治疗时的风险：$\Pr[Y=1\given L=0,A=1]=\Pr[Y^{a=1}=1\given L=0]$；未受治疗者的风险等于所有人都未受治疗时的风险：$\Pr[Y=1\given L=0,A=0]=\Pr[Y^{a=0}=1\given L=0]$。同理可断定，12 名危重者一组中观测风险也等于反事实风险，即 $\Pr[Y=1\given L=1,A=1]=\Pr[Y^{a=1}=1\given L=1]=2/3$，$\Pr[Y=1\given L=1,A=0]=\Pr[Y^{a=0}=1\given L=1]=2/3$。

现在设我们的目标是计算因果风险比 $\Pr[Y^{a=1}=1]/\Pr[Y^{a=0}=1]$。因果风险比的分子，是总体中全部 20 人都受治疗时的风险。由上一段可知，所有人都受治疗时的风险在 $L=0$ 的 8 人中为 $1/4$，在 $L=1$ 的 12 人中为 $2/3$。因此，全部 20 人都受治疗时的风险，是 $1/4$ 与 $2/3$ 的加权平均，各组的权重与其规模成比例。由于 40\% 的人（8 人）属于 $L=0$ 组、60\% 的人（12 人）属于 $L=1$ 组，加权平均为 $\frac{1}{4}\times 0.4+\frac{2}{3}\times 0.6=0.5$。于是所有人都受治疗时的风险 $\Pr[Y^{a=1}=1]$ 等于 0.5。同理可算得，所有人都不受治疗时的风险 $\Pr[Y^{a=0}=1]$ 也等于 0.5。因果风险比于是为 $0.5/0.5=1$。

更正式地说，边际反事实风险 $\Pr[Y^{a}=1]$ 是层特异风险 $\Pr[Y^{a}=1\given L=0]$ 与 $\Pr[Y^{a}=1\given L=1]$ 的加权平均，权重分别等于总体中 $L=0$ 与 $L=1$ 的个体所占比例。即
\[
\Pr[Y^{a}=1]=\Pr[Y^{a}=1\given L=0]\Pr[L=0]+\Pr[Y^{a}=1\given L=1]\Pr[L=1].
\]
或者用更紧凑的记号，
\[
\Pr[Y^{a}=1]=\sum_{l}\Pr[Y^{a}=1\given L=l]\Pr[L=l],
\]
其中 $\sum_l$ 表示对总体中出现的一切取值 $l$ 求和。在条件可交换性之下，我们可以把上式中的反事实风险 $\Pr[Y^{a}=1\given L=l]$ 换成观测风险 $\Pr[Y=1\given L=l,A=a]$，即
\[
\Pr[Y^{a}=1]=\sum_{l}\Pr[Y=1\given L=l,A=a]\Pr[L=l].
\]
这一等式的左端是未观测到的反事实风险，右端只含观测量，可用 $L$、$A$、$Y$ 的数据算出。当反事实量能像这里这样表示为观测数据分布（即概率）的函数时，我们说该反事实量是\emph{可识别的}；否则称之为\emph{不可识别的}。

\mnote{标准化均值 $\sum_l\E[Y\given L=l,A=a]\times\Pr[L=l]$}

这种方法在流行病学、人口学及其他学科中称为\emph{标准化}。例如，因果风险比的分子 $\sum_l\Pr[Y=1\given L=l,A=1]\Pr[L=l]$，就是以总体为标准人群的受治疗者的标准化风险。在条件可交换性之下，这一标准化风险可以解释为：假如总体中所有人都受治疗，本会观测到的（反事实）风险。受治疗者与未受治疗者的标准化风险，分别等于治疗之下与未治疗之下的反事实风险。因此，因果风险比 $\Pr[Y^{a=1}=1]/\Pr[Y^{a=0}=1]$ 可用标准化计算如下：
\[
\frac{\sum_l\Pr[Y=1\given L=l,A=1]\Pr[L=l]}{\sum_l\Pr[Y=1\given L=l,A=0]\Pr[L=l]}.
\]

\section{逆概率加权}

上一节我们用标准化计算了条件随机化实验中的因果风险比。这一节改用逆概率加权计算同一因果风险比。表~2.2 的数据可以画成一棵树：全部 20 人从左侧出发，随时间向右推进，如图~2.1。树最左边的圆圈包含第一次分叉：8 人病情不危重（$L=0$），12 人病情危重（$L=1$）。括号中的数是处于不危重状态的概率 $\Pr[L=0]=8/20=0.4$ 与危重状态的概率 $\Pr[L=1]=12/20=0.6$。以 $L=0$ 这一支为例：这 8 人中有 4 人未受治疗（$A=0$）、4 人受治疗（$A=1$）。未受治疗的条件概率为 $\Pr[A=0\given L=0]=4/8=0.5$，如括号所示；受治疗的条件概率 $\Pr[A=1\given L=0]$ 也是 0.5。右上方的圆圈表示：$(L=0,A=0)$ 这一支的 4 人中，3 人存活（$Y=0$）、1 人死亡（$Y=1$），即 $\Pr[Y=0\given L=0,A=0]=3/4$，$\Pr[Y=1\given L=0,A=0]=1/4$。树的其余各支依此类推。圆圈中包含的是由非治疗变量定义的分叉。下面用这棵树计算因果风险比。

\mnote{图 2.1 是一个“完全随机化因果解释结构树图”（fully randomized causally interpreted structured tree graph, FR-CISTG，Robins 1986, 1987）的例子，用来表示条件随机化实验。我们是不是该得一个“史上最糟缩写”奖？}

\begin{figure}[t]
\centering
\begin{tikzpicture}[causal, every node/.style={transform shape}]
  \draw[thin] (0,0) ellipse (0.95 and 0.55);
  \draw[thin] (-0.86,-0.14) -- (0.86,0.14);
  \node[font=\footnotesize,inner sep=0pt] at (0,0.22) {$L{=}0$ (.4)};
  \node[font=\footnotesize,inner sep=0pt] at (0,-0.24) {$L{=}1$ (.6)};
  \ciLeafA{6.6}{3.0}{$Y{=}0$ (.75)}{$Y{=}1$ (.25)}{3}{1}
  \ciLeafA{6.6}{1.5}{$Y{=}0$ (.75)}{$Y{=}1$ (.25)}{3}{1}
  \ciLeafA{6.6}{-1.5}{$Y{=}0$ (1/3)}{$Y{=}1$ (2/3)}{1}{2}
  \ciLeafA{6.6}{-3.0}{$Y{=}0$ (1/3)}{$Y{=}1$ (2/3)}{3}{6}
  \draw[thin] (0.95,0) -- node[pos=0.52,above,sloped,font=\footnotesize,inner sep=1pt] {$A{=}0$ (.5)}
      node[pos=0.80,below,sloped,font=\footnotesize,inner sep=1pt] {4} (5.30,3.0);
  \draw[thin] (0.95,0) -- node[pos=0.52,above,sloped,font=\footnotesize,inner sep=1pt] {$A{=}1$ (.5)}
      node[pos=0.80,below,sloped,font=\footnotesize,inner sep=1pt] {4} (5.30,1.5);
  \draw[thin] (0.95,0) -- node[pos=0.52,above,sloped,font=\footnotesize,inner sep=1pt] {$A{=}0$ (.25)}
      node[pos=0.80,below,sloped,font=\footnotesize,inner sep=1pt] {3} (5.30,-1.5);
  \draw[thin] (0.95,0) -- node[pos=0.52,above,sloped,font=\footnotesize,inner sep=1pt] {$A{=}1$ (.75)}
      node[pos=0.80,below,sloped,font=\footnotesize,inner sep=1pt] {9} (5.30,-3.0);
\end{tikzpicture}
\caption{表 2.2 的数据构成的树：圆圈内为分叉的概率，圆圈右侧为人数}
\label{fig:2.1}
\end{figure}

因果风险比的分母 $\Pr[Y^{a=0}=1]$，是总体中所有人都未受治疗时的反事实死亡风险。我们来算这个风险。图~2.1 中，$L=0$ 的 8 人里有 4 人未受治疗，其中 1 人死亡。假如 $L=0$ 的 8 人全都未受治疗，本会发生多少例死亡？两例。因为若未受治疗的是 8 人而不是 4 人，观测到的死亡就会是 2 例而不是 1 例。人数乘以 2，死亡数也加倍。图~2.1 中，$L=1$ 的 12 人里有 3 人未受治疗，其中 2 人死亡。假如 $L=1$ 的 12 人全都未受治疗，本会发生多少例死亡？八例，即 2 例乘以 4，因为 12 是 $3\times 4$。也就是说，若总体中全部 $8+12=20$ 人都未受治疗，本会有 $2+8=10$ 人死亡。因果风险比的分母 $\Pr[Y^{a=0}=1]$ 为 $10/20=0.5$。图~2.2 中的第一棵树画出了所有人都未受治疗时的总体。当然，这些计算依赖一个条件：$L=0$ 的受治疗者若一直未受治疗，其死亡概率本会与真正未受治疗者相同。这个条件正是给定 $L=0$ 的可交换性。

因果风险比的分子 $\Pr[Y^{a=1}=1]$，是总体中所有人都受治疗时的反事实死亡风险。按上一段的思路，在给定 $L=1$ 的可交换性之下，算得这一风险也是 $10/20=0.5$。图~2.2 中的第二棵树画出了所有人都受治疗时的总体。综合本段与上一段的结果，因果风险比 $\Pr[Y^{a=1}=1]/\Pr[Y^{a=0}=1]$ 等于 $0.5/0.5=1$。到此完成。

\begin{figure}[t]
\centering
\begin{tikzpicture}[causal, every node/.style={transform shape}]
  % 上树：所有人都未受治疗
  \begin{scope}
    \draw[thin] (0,0) ellipse (0.80 and 0.44);
    \draw[thin] (-0.73,-0.11) -- (0.73,0.11);
    \node[font=\footnotesize,inner sep=0pt] at (0,0.18) {$L{=}0$ (.4)};
    \node[font=\footnotesize,inner sep=0pt] at (0,-0.20) {$L{=}1$ (.6)};
    \ciLeafB{6.0}{2.2}{$Y{=}0$ (.75)}{$Y{=}1$ (.25)}{6}{2}
    \ciLeafB{6.0}{0.7}{$Y{=}0$}{$Y{=}1$}{0}{0}
    \ciLeafB{6.0}{-0.7}{$Y{=}0$ (1/3)}{$Y{=}1$ (2/3)}{4}{8}
    \ciLeafB{6.0}{-2.2}{$Y{=}0$}{$Y{=}1$}{0}{0}
    \draw[thin] (0.80,0) -- node[pos=0.50,above,sloped,font=\scriptsize,inner sep=1pt] {Untreated(1)}
        node[pos=0.72,below,sloped,font=\scriptsize,inner sep=1pt] {8} (5.00,2.2);
    \draw[thin] (0.80,0) -- node[pos=0.64,above,sloped,font=\scriptsize,inner sep=1pt] {Treated(0)}
        node[pos=0.84,below,sloped,font=\scriptsize,inner sep=1pt] {0} (5.00,0.7);
    \draw[thin] (0.80,0) -- node[pos=0.64,above,sloped,font=\scriptsize,inner sep=1pt] {Untreated(1)}
        node[pos=0.84,below,sloped,font=\scriptsize,inner sep=1pt] {12} (5.00,-0.7);
    \draw[thin] (0.80,0) -- node[pos=0.50,above,sloped,font=\scriptsize,inner sep=1pt] {Treated(0)}
        node[pos=0.72,below,sloped,font=\scriptsize,inner sep=1pt] {0} (5.00,-2.2);
  \end{scope}
  % 下树：所有人都受治疗
  \begin{scope}[shift={(0,-5.6)}]
    \draw[thin] (0,0) ellipse (0.80 and 0.44);
    \draw[thin] (-0.73,-0.11) -- (0.73,0.11);
    \node[font=\footnotesize,inner sep=0pt] at (0,0.18) {$L{=}0$ (.4)};
    \node[font=\footnotesize,inner sep=0pt] at (0,-0.20) {$L{=}1$ (.6)};
    \ciLeafB{6.0}{2.2}{$Y{=}0$}{$Y{=}1$}{0}{0}
    \ciLeafB{6.0}{0.7}{$Y{=}0$ (.75)}{$Y{=}1$ (.25)}{6}{2}
    \ciLeafB{6.0}{-0.7}{$Y{=}0$}{$Y{=}1$}{0}{0}
    \ciLeafB{6.0}{-2.2}{$Y{=}0$ (1/3)}{$Y{=}1$ (2/3)}{4}{8}
    \draw[thin] (0.80,0) -- node[pos=0.50,above,sloped,font=\scriptsize,inner sep=1pt] {Untreated(0)}
        node[pos=0.72,below,sloped,font=\scriptsize,inner sep=1pt] {0} (5.00,2.2);
    \draw[thin] (0.80,0) -- node[pos=0.64,above,sloped,font=\scriptsize,inner sep=1pt] {Treated(1)}
        node[pos=0.84,below,sloped,font=\scriptsize,inner sep=1pt] {8} (5.00,0.7);
    \draw[thin] (0.80,0) -- node[pos=0.64,above,sloped,font=\scriptsize,inner sep=1pt] {Untreated(0)}
        node[pos=0.84,below,sloped,font=\scriptsize,inner sep=1pt] {0} (5.00,-0.7);
    \draw[thin] (0.80,0) -- node[pos=0.50,above,sloped,font=\scriptsize,inner sep=1pt] {Treated(1)}
        node[pos=0.72,below,sloped,font=\scriptsize,inner sep=1pt] {12} (5.00,-2.2);
  \end{scope}
\end{tikzpicture}
\caption{所有人都未受治疗（上）与所有人都受治疗（下）时的总体。括号内为条件概率，圆圈右侧为人数}
\label{fig:2.2}
\end{figure}

来看这一方法为何有效。图~2.2 的两棵树是对以下两种情形的模拟：总体中所有人都未受治疗，以及所有人都受治疗。在条件可交换性之下，这两次模拟都是正确的。把两次模拟合并，就得到一个假想总体：其中每个人都以受治疗者与未受治疗者两种身份各出现一次。这个假想总体的规模是原总体的两倍，称为\emph{伪总体}。图~2.3 画出了整个伪总体。在原总体中条件可交换性 $Y^{a}\indep A\given L$ 成立时，伪总体中的受治疗者与未受治疗者（无条件地）可交换，因为 $L$ 与 $A$ 独立。也就是说，伪总体中的关联风险比等于伪总体与原总体中的因果风险比。

\begin{figure}[t]
\centering
\begin{tikzpicture}[causal, every node/.style={transform shape}]
  \draw[thin] (0,0) ellipse (0.95 and 0.55);
  \draw[thin] (-0.86,-0.14) -- (0.86,0.14);
  \node[font=\footnotesize,inner sep=0pt] at (0,0.22) {$L{=}0$ (.4)};
  \node[font=\footnotesize,inner sep=0pt] at (0,-0.24) {$L{=}1$ (.6)};
  \ciLeafA{6.6}{3.0}{$Y{=}0$ (.75)}{$Y{=}1$ (.25)}{6}{2}
  \ciLeafA{6.6}{1.5}{$Y{=}0$ (.75)}{$Y{=}1$ (.25)}{6}{2}
  \ciLeafA{6.6}{-1.5}{$Y{=}0$ (1/3)}{$Y{=}1$ (2/3)}{4}{8}
  \ciLeafA{6.6}{-3.0}{$Y{=}0$ (1/3)}{$Y{=}1$ (2/3)}{4}{8}
  \draw[thin] (0.95,0) -- node[pos=0.52,above,sloped,font=\footnotesize,inner sep=1pt] {Untreated(1)}
      node[pos=0.80,below,sloped,font=\footnotesize,inner sep=1pt] {8} (5.30,3.0);
  \draw[thin] (0.95,0) -- node[pos=0.52,above,sloped,font=\footnotesize,inner sep=1pt] {Treated(1)}
      node[pos=0.80,below,sloped,font=\footnotesize,inner sep=1pt] {8} (5.30,1.5);
  \draw[thin] (0.95,0) -- node[pos=0.52,above,sloped,font=\footnotesize,inner sep=1pt] {Untreated(1)}
      node[pos=0.80,below,sloped,font=\footnotesize,inner sep=1pt] {12} (5.30,-1.5);
  \draw[thin] (0.95,0) -- node[pos=0.52,above,sloped,font=\footnotesize,inner sep=1pt] {Treated(1)}
      node[pos=0.80,below,sloped,font=\footnotesize,inner sep=1pt] {12} (5.30,-3.0);
  \node[font=\footnotesize,anchor=west,inner sep=0pt] at (8.95,3.72) {$W^{A}=1/f[A\given L]$};
  \node[font=\footnotesize,anchor=west,inner sep=0pt] at (8.95,3.31) {$1/.5{=}2$};
  \node[font=\footnotesize,anchor=west,inner sep=0pt] at (8.95,2.69) {$1/.5{=}2$};
  \node[font=\footnotesize,anchor=west,inner sep=0pt] at (8.95,1.81) {$1/.5{=}2$};
  \node[font=\footnotesize,anchor=west,inner sep=0pt] at (8.95,1.19) {$1/.5{=}2$};
  \node[font=\footnotesize,anchor=west,inner sep=0pt] at (8.95,-1.19) {$1/.25{=}4$};
  \node[font=\footnotesize,anchor=west,inner sep=0pt] at (8.95,-1.81) {$1/.25{=}4$};
  \node[font=\footnotesize,anchor=west,inner sep=0pt] at (8.95,-2.69) {$1/.75{=}1.33$};
  \node[font=\footnotesize,anchor=west,inner sep=0pt] at (8.95,-3.31) {$1/.75{=}1.33$};
\end{tikzpicture}
\caption{逆概率加权构造的伪总体：最右列为各人的逆概率权重 $W^{A}$}
\label{fig:2.3}
\end{figure}

这种方法称为\emph{逆概率}（inverse probability, IP）\emph{加权}。要看清缘由，不妨看图~2.1 总体中 $L=0$ 的 4 名未受治疗者。这 4 人被用来构造图~2.3 伪总体中的 8 个成员，也就是说，每人都分到一个权重 2，正好等于 $1/0.5$。图~2.1 显示，0.5 是给定 $L=0$ 时保持未受治疗的条件概率。同样，图~2.1 中 $L=1$ 的 9 名受治疗者被用来构造伪总体中的 12 个成员，每人的权重为 $1.33=1/0.75$。\mnote{逆概率权重：$W^{A}=1/f[A\given L]$}图~2.1 显示，0.75 是给定 $L=1$ 时受治疗的条件概率。通俗地说，构造伪总体的办法，是把总体中每个人按“接受其实际所接受的那个治疗水平的条件概率”的倒数加权。这些逆概率权重见图~2.3。

\mnote{逆概率加权估计量由 Horvitz 与 Thompson (1952) 提出，用于个体以不等概率被抽中的抽样调查。见技术注记 12.1。}

\begin{techpoint}[逆概率权重的形式定义]
一个个体的逆概率权重取决于该个体的治疗取值 $A$ 与协变量 $L$。例如，$L=l$ 的受治疗者分到权重 $1/\Pr[A=1\given L=l]$，而 $L=l'$ 的未受治疗者分到权重 $1/\Pr[A=0\given L=l']$。要用一个式子表示所有个体的权重——不论其治疗与协变量取值如何——我们可以使用 $A$ 的概率密度函数（pdf），而不用 $A$ 的概率。给定 $L$ 时 $A$ 的条件密度在取值 $a$ 与 $l$ 处的值记为 $f_{A|L}[a|l]$，或简记为 $f[a|l]$。对离散变量 $A$ 与 $L$，$f[a|l]$ 就是条件概率 $\Pr[A=a\given L=l]$。在条件随机化实验中，凡 $\Pr[L=l]$ 非零的 $l$，$f[a|l]$ 都为正。

由于每个个体权重的分母，是在该个体自身的 $A$ 与 $L$ 取值处取值的条件密度，它可以写成以随机变量 $A$ 与 $L$（而非固定的 $a$ 与 $l$）为自变量的条件密度，即 $f[A|L]$。图~2.3 中出现的就是这个记号，逆概率权重即定义为 $W^{A}=1/f[A|L]$。这里需要一种统一的权重记号，因为 $\Pr[A=A\given L=L]$ 同义反复地等于 1，不能算作恰当的记号。

如正文所述，伪总体中结局的均值 $\E_{ps}[Y\given A=a]$ 等于总体中结局的逆概率加权均值 $\E[YI(A=a)/\Pr(A=a\given L)]$，其中 $I(A=a)$ 在 $A=a$ 时为 1、否则为 0。证明如下：
\begin{align*}
\E_{ps}[Y\given A=a] &= \frac{\E_{ps}[YI(A=a)]}{\E_{ps}[I(A=a)]}\\
&= \frac{\E\bigl[W^{A}YI(A=a)\bigr]}{\E\bigl[I(A=a)W^{A}\bigr]}\\
&= \frac{\E\bigl[YI(A=a)/\Pr(a\given L)\bigr]}{\E\bigl[I(A=a)/\Pr(a\given L)\bigr]}\\
&= \E\bigl[YI(A=a)/\Pr(a\given L)\bigr].
\end{align*}
其中第一行由概率法则；第二行由伪总体均值 $\E_{ps}$ 的定义；第三行因 $I(A=a)/f[A|L]=I(A=a)/f[a|L]$，而后者等于 $I(A=a)/\Pr(a\given L)$；第四行因 $\E[I(A=a)/\Pr(a\given L)\given L]=1$。
\end{techpoint}

在上面的例子里，逆概率加权与标准化给出了相同的结果——因果风险比都等于 1。这并非巧合：标准化与逆概率加权在数学上等价（见技术注记 2.3）。事实上，标准化与逆概率加权都可视为构造一棵新树的程序，在这棵新树中所有人都接受治疗 $a$。两种方法用来构造反事实树的那组概率不同：逆概率加权用给定协变量 $L$ 时治疗 $A$ 的条件概率（如图~2.1 所示），标准化用协变量 $L$ 的概率以及给定 $A$ 与 $L$ 时结局 $Y$ 的条件概率。

标准化与逆概率加权都在模拟这样一种情形：假如没有用变量（或向量中的诸变量）$L$ 来决定治疗概率，本会观测到什么。因此我们常说，这些方法对 $L$ 作了\emph{调整}。说得随便一点，我们有时说这些方法\emph{控制}了 $L$，但这种“分析上的控制”与随机化实验中的“物理控制”大不相同。标准化与逆概率加权可以推广到结局为连续变量的条件随机化研究（见技术注记 2.3）。

本书何不就此收尾？我们已经有了一种研究设计（理想的随机化实验），配合恰当的分析方法（标准化或逆概率加权），就能算出平均因果效应。遗憾的是，随机化实验往往不合伦理、不切实际或赶不上时机。例如，伦理委员会是否会批准我们这项心脏移植研究，就很成问题。心脏供不应求，社会倾向于把它们分配给最可能从移植中获益的人，而不是在潜在受者之间随机分配。即便不考虑伦理问题，研究的可行性也值得怀疑：双盲分配不可能做到；被分配到药物治疗的人未必甘心放弃移植；被分配到移植的人未必都有配型相合的心脏。就算研究可行，完成它也要好几年，而决策必须在这期间作出。经常的情况是，开展观察性研究已是坏处最少的选择。

\begin{techpoint}[逆概率加权与标准化的等价性]
假定 $A$ 为离散变量、取值个数有限，且凡 $\Pr[L=l]$ 非零的 $l$ 都有 $f[a|l]>0$。这一\emph{正值性}条件在条件随机化实验中必然成立。在正值性之下，治疗水平 $a$ 的标准化均值定义为
\[
\sum_l\E[Y\given A=a,L=l]\Pr[L=l],
\]
而 $Y$ 在治疗水平 $a$ 的逆概率加权均值定义为
\[
\E\left[\frac{I(A=a)Y}{f[A|L]}\right].
\]
指示函数 $I(A=a)$ 对 $A=a$ 的个体取 1、对其余个体取 0。

现在证明在正值性之下逆概率加权均值与标准化均值相等。由期望的定义，
\begin{align*}
\E\left[\frac{I(A=a)Y}{f[A|L]}\right]
&=\sum_l\frac{1}{f[a|l]}\Big\{\E[Y\given A=a,L=l]\,f[a|l]\Pr[L=l]\Big\}\\
&=\sum_l\Big\{\E[Y\given A=a,L=l]\Pr[L=l]\Big\},
\end{align*}
其中最后一步从分子与分母中约去了 $f[a|l]$；第一步无需对 $A$ 的可能取值求和，因为对任何不等于 $a$ 的 $a'$，量 $I(a'=a)$ 都为零。证明把 $A$ 与 $L$ 当作离散变量处理，但不要求二值；对连续的 $L$，只需把对 $L$ 的求和换成积分。

这一证明没有用到反事实。然而，若进一步假定条件可交换性，则逆概率加权均值与标准化均值都等于反事实均值 $\E[Y^{a}]$。下面给出这最后一个论断的两种不同证明。第一种，如正文那样证明 $\E[Y^{a}]$ 与标准化均值相等：
\begin{align*}
\E[Y^{a}]&=\sum_l\E[Y^{a}\given L=l]\Pr[L=l]\\
&=\sum_l\E[Y^{a}\given A=a,L=l]\Pr[L=l]\\
&=\sum_l\E[Y\given A=a,L=l]\Pr[L=l],
\end{align*}
其中第二个等号由条件可交换性与正值性得到，第三个等号由一致性得到。第二种，证明 $\E[Y^{a}]$ 与逆概率加权均值相等。首先，由一致性有
\[
\E\left[\frac{I(A=a)}{f[A|L]}Y^{a}\right]=\E\left[\frac{I(A=a)}{f[A|L]}Y\right].
\]
其次，由于正值性蕴含 $f[a|L]$ 永不为 0，我们有
\begin{align*}
\E\left[\frac{I(A=a)}{f[a|L]}Y^{a}\right]
&=\E\left\{\E\left[\frac{I(A=a)}{f[a|L]}Y^{a}\Big|L\right]\right\}\\
&=\E\left\{\E\left[\frac{I(A=a)}{f[a|L]}\Big|L\right]\E[Y^{a}|L]\right\}\\
&=\E\{\E[Y^{a}|L]\}=\E[Y^{a}],
\end{align*}
其中第二个等号由条件可交换性得到，第三个等号因 $\E[I(A=a)/f[a|L]\given L]=1$。

当治疗为连续变量时——这在条件随机化实验中不太可能是设计选择——$\E[I(A=a)Y/f(A|L)]$ 不再等于 $\sum_l\E[Y\given A=a,L=l]\Pr[L=l]$，因而即便在可交换性之下也是有偏的。要看这一点，可以算出：若把 $f(a|l)$ 取为给定 $L=l$ 时 $A$ 的条件密度（关于 Lebesgue 测度的某个版本），则 $\E[I(A=a)/f(a|l)\given L=l]$ 等于 0 而不是 1。另一方面，若继续把 $f(a|l)$ 取为 $\Pr[A=a\given L=l]$，则分母 $f(a\given L=l)$ 在一个概率为 1 的集合上为零，正值性不成立。12.4 节讨论如何把逆概率加权推广到连续治疗。技术注记 3.1 讨论：没有正值性时，上述结果不成立，即便 $A$ 是离散的。
\end{techpoint}
